% Part: many-valued-logic
% Chapter: syntax-and-semantics
% Section: matrices

\documentclass[../../../include/open-logic-section]{subfiles}

\begin{document}

\olfileid{mvl}{syn}{mat}

\olsection{શ્રેણિકો}

બહુમૂલ્ય તર્કશાસ્ત્ર તેની ભાષા, સત્યમૂલ્યોના ગણ~$V$,
નિર્દિષ્ટ સત્યમૂલ્યોના ઉપગણ અને તેના સંયોજકોનાં
સત્યવિધેયો વડે વ્યાખ્યાયિત થાય છે. આ ઘટકોને
એકસાથે \emph{શ્રેણિક} કહે છે.

\begin{defn}[શ્રેણિક]
\ollabel{defn:matrix}
તર્કશાસ્ત્ર~$\Log L$ માટેની \emph{શ્રેણિક}માં આ હોય છે:
\begin{enumerate}
\item ભાષા~$\Lang L$ બનાવતો સંયોજકોનો ગણ;
\item સત્યમૂલ્યોનો અરિક્ત ગણ~$V \neq \emptyset$;
\item નિર્દિષ્ટ સત્યમૂલ્યોનો ગણ~$V^+ \subseteq V$;
\item $\Lang L$ના દરેક $n$-સ્થાનીય સંયોજક~$\star$
માટે સત્યવિધેય~$\tf{\star} : V^n \to V$. જો
$n = 0$ હોય, તો $\tf{\star}$ એ $V$નો માત્ર એક ઘટક છે.
\end{enumerate}
\end{defn}

\begin{ex}
શાસ્ત્રીય તર્કશાસ્ત્ર~\LogCL{} માટેની શ્રેણિકમાં આ હોય છે:
\begin{enumerate}
  \item $\lfalse$, $\lnot$, $\land$, $\lor$ અને $\lif$
  ધરાવતી પ્રમાણભૂત વિધાનાત્મક ભાષા~$\Lang L_0$.
  \item સત્યમૂલ્યોનો ગણ~$V = \{\True, \False\}$.
  \item એકમાત્ર નિર્દિષ્ટ મૂલ્ય~$\True$ છે, એટલે
  $V^+ = \{\True\}$.
  \item $\lfalse$ માટે $\tf{\lfalse} = \False$ છે.
  અન્ય સત્યવિધેયો સામાન્ય સત્યતાસારણીઓ વડે
  અપાય છે (જુઓ \olref{fig:tf-CL}).
\end{enumerate}
\begin{figure}
  \begin{center}
      \begin{tabular}{c|c}
        $\tf{\lnot}$ & \\
        \hline
        $\True$ & $\False$ \\
        $\False$ & $\True$
      \end{tabular}
      \quad
      \begin{tabular}{c|cc}
        $\tf{\land}$ & $\True$ & $\False$ \\
        \hline
        $\True$ & $\True$ & $\False$ \\
        $\False$ & $\False$ & $\False$
      \end{tabular}
      \quad
      \begin{tabular}{c|cc}
        $\tf{\lor}$ & $\True$ & $\False$ \\
        \hline
        $\True$ & $\True$ & $\True$ \\
        $\False$ & $\True$ & $\False$
      \end{tabular}
      \quad
      \begin{tabular}{c|cc}
        $\tf{\lif}$ & $\True$ & $\False$ \\
        \hline
        $\True$ & $\True$ & $\False$ \\
        $\False$ & $\True$ & $\True$
      \end{tabular}
    \end{center}
    \caption{શાસ્ત્રીય તર્કશાસ્ત્ર~$\LogCL$ માટેનાં સત્યવિધેયો.}
    \ollabel{fig:tf-CL}
  \end{figure}
  \end{ex}

\end{document}
